\documentclass[11pt]{article}
\usepackage[margin=.85in,headheight=15pt]{geometry}
\usepackage[T1]{fontenc}
\usepackage{lmodern,microtype,mathtools,amssymb,booktabs,xcolor}
\usepackage{hyperref,xurl,fancyhdr,needspace}
\definecolor{ink}{HTML}{17324D}
\definecolor{teal}{HTML}{167D82}
\color{ink}
\hypersetup{colorlinks=true,urlcolor=teal,linkcolor=teal,
pdftitle={Presenter notes: finite fields and Fibonacci, four 50-minute lectures},
pdfauthor={Igor Rivin}}
\setlength{\parindent}{0pt}
\setlength{\parskip}{5pt}
\pagestyle{fancy}
\fancyhf{}
\fancyhead[L]{\small Finite fields and Fibonacci: presenter notes}
\fancyhead[R]{\small Four 50-minute lectures}
\fancyfoot[C]{\thepage}
\newcommand{\Q}{\mathbb Q}
\newcommand{\Z}{\mathbb Z}
\newcommand{\R}{\mathbb R}
\newcommand{\F}{\mathbb F}
\newcommand{\GL}{\operatorname{GL}}
\newcommand{\SL}{\operatorname{SL}}
\newcommand{\Span}{\operatorname{span}}
\newcommand{\ord}{\operatorname{ord}}
\newcommand{\mat}[4]{\begin{pmatrix}#1&#2\\#3&#4\end{pmatrix}}
\newcounter{slide}
\newcounter{elapsed}
\newcounter{backup}
\newcommand{\teachslide}[4]{%
  \stepcounter{slide}\Needspace{8\baselineskip}
  \subsection*{\theslide.\ #1}
  {\small\color{teal}\textbf{Minutes \theelapsed\ to
  \number\numexpr\value{elapsed}+#2\relax}\par}
  \addtocounter{elapsed}{#2}#4\par}
\let\teachcover\teachslide
\newcommand{\backups}{%
  \typeout{LECTURE DURATION: \theelapsed\space minutes; CORE SLIDES: \theslide}
  \ifnum\value{elapsed}=50\else\errmessage{Lecture timing is not 50 minutes}\fi
  \setcounter{backup}{0}
  \Needspace{8\baselineskip}\section*{Optional slides after the 50-minute stopping point}}
\newcommand{\backupslide}[3]{%
  \stepcounter{backup}\Needspace{8\baselineskip}
  \subsection*{Backup \thebackup.\ #1}#3\par}
\begin{document}
{\huge\bfseries Finite fields and Fibonacci\par}
{\Large Presenter notes for four 50-minute lectures\par}
\vspace{10pt}
These lectures begin with the matrix method for Fibonacci periodicity,
develop rings, fields, characteristic, and vector spaces, prove cyclicity
of finite-field multiplication, and return to the prime-modulus period
through matrix groups and a quadratic extension.
The visual style follows the earlier continued-fraction lectures.

\textbf{Before Lecture 1.}
Students should know modular arithmetic, Bezout's identity, modular
inverses, $2\times2$ matrix multiplication and determinants, and the
Fibonacci recurrence. The slides review their relevant consequences.
The vector-space reminder assumes only some prior contact with linear
algebra and develops span, independence, bases, coordinates, and dimension.
Lecture 2 proves that dimension is well defined by the exchange argument
before using dimension to count elements. Gaussian elimination is an
optional alternative proof, not a prerequisite.

\textbf{Before Lecture 3.}
Recall Lagrange's theorem for finite groups and Bezout's identity.
The deck proves the polynomial root bound and the coprime-order product
lemma used to construct a generator. No classification theorem for
finite abelian groups or general Galois theory is used.

\textbf{Before Lecture 4.}
Recall Fermat's little theorem and divisibility of the intermediate
binomial coefficients $\binom pj$ by a prime $p$. The deck supplies
explicit eigenvectors and proves the Frobenius argument, using Lecture 3's
polynomial root bound to identify the fixed elements.

\textbf{Conventions.}
Our rings are commutative with identity and $1\ne0$. A field has
characteristic $0$ or a prime; a finite field necessarily has prime
characteristic. We use $d$ for a vector-space dimension to avoid
confusing it with an integer modulus $n$. We use
$A=\begin{psmallmatrix}1&1\\1&0\end{psmallmatrix}$ and order the matrix
state as $(F_{n+1},F_n)^{\mathsf T}$. A restart is still described by
$(F_n,F_{n+1})=(0,1)$. For odd $p$, $A$ has determinant $-1$ and
$A^2$ belongs to $\SL_2(\F_p)$.

\textbf{Virtual delivery.}
Share the slide PDF window and keep these notes separate. Each question
page precedes its answer page, so advancing the PDF reveals the answer.
Give students quiet calculation time before inviting chat or spoken
responses. The suggested minute allocations include that work. No
animations, browser applications, or live network access are needed.

\textbf{Pacing and recovery.}
Each main route totals 50 minutes, including its cover and exit question.
Lecture 1 reaches vector spaces after 32 minutes. Lecture 2 reaches the
field construction after 36 minutes. Its main slides 7--12 spend minutes
12--24 on well-definedness of dimension, including a worked exchange.
Lecture 3 completes cyclicity at minute 28 and uses the next 20
minutes to count GL and SL; the exit question finishes the session.
Lecture 4 reaches Frobenius at minute 29, leaving room for its proof
and both period examples.
These are pacing guides, not deadlines for students. If Lecture 1 runs
long, shorten the recap and answer reveals, keeping student work. If
Lecture 2 runs long, shorten the opening recap or the repeated pair
example, keeping the exchange proof and the basis exercise. The
extension-arithmetic exercise is already optional; the main route
uses a worked multiplication and retains the inverse proof.
In Lecture 3, the coprime-order lemma and
polynomial argument are the essential new proofs; omit the general
matrix count or the modulo-two exercise if students need more time.
In Lecture 4, shorten the recap and the primitive-element comparison
before shortening the Frobenius argument.

\textbf{The optional prime-power connection.}
Lecture 3's extra slides compare $\F_{p^n}$ with $\Z/p^n\Z$. The latter
is a different object: for $n>1$ it has zero divisors, yet its unit group
is cyclic when $p$ is odd. The slides give a binomial primitive-root
lifting proof and state the exact exception for powers of $2$.
Unit groups modulo $2$ and $4$ are cyclic; modulo $2^n$ for $n\ge3$ they
are $C_2\times C_{2^{n-2}}$. The notes justify that classification using
$-1$ and $5$. Budget another 15--20 minutes if teaching the entire
optional proof; it is outside the four 50-minute main routes.

\section*{Lecture 1: Changing the arithmetic}
\teachcover{Changing the arithmetic}{1}{
Rings, fields, and vectors
\[
X^2-X-1=0\qquad\text{modulo }p
\]
}{
The students already know Fibonacci periodicity via a matrix. Today names
the arithmetic they have been using and begins the vector-space reminder.
The goal is not diagonalization yet. Allow calculation time at slides 8
and 20. All rings in these lectures are commutative with identity, with
$1\ne0$. A field has the additional inverse property.
}

\teachslide{The matrix we already know}{2}{
\[
A=\mat1110,\qquad
\begin{pmatrix}F_{n+1}\\F_n\end{pmatrix}
=A^n\begin{pmatrix}1\\0\end{pmatrix}.
\]
\[
\det A=-1,\qquad A^2=A+I.
\]
Modulo any integer $m\ge2$, the step is reversible.

\key{Today: what arithmetic are these entries using?}
}{
Use $F_0=0,F_1=1$. This matrix orders the state as $(F_{n+1},F_n)$;
the book sometimes uses the reversed state and the conjugate matrix.
They give the same period. Ask students to multiply $A$ by itself.
The determinant is a unit modulo every modulus, so the inverse exists.
Do not confuse invertibility of this particular matrix with every
nonzero residue being invertible.
}

\teachslide{Where are the roots?}{2}{
\[
\det(XI-A)=X^2-X-1.
\]
Over the real numbers, its roots are
\[
\frac{1+\sqrt5}{2},\qquad \frac{1-\sqrt5}{2}.
\]
\key{Modulo $11$:}\quad $4^2=5$.

\key{Modulo $7$:}\quad squares are $0,1,2,4$.
}{
Reduce $16$ modulo $11$. For $7$, list the squares of $0,1,2,3$;
the remaining inputs give the same squares. Stress that writing a
square-root symbol does not manufacture a residue with that property.
Division by $2$ will make sense for odd primes. Prime $2$ is exceptional
and can wait until Lecture 3's optional material.
}

\teachslide{Residues are numbers with arithmetic}{2}{
\[
\Z/n\Z=\{[0],[1],\ldots,[n-1]\}.
\]
\[
[a]+[b]=[a+b],\qquad [a][b]=[ab].
\]
For example, in $\Z/6\Z$,
\[
[4]+[5]=[3],\qquad [4][5]=[2].
\]
We usually omit the brackets once the modulus is fixed.
}{
The objects are residue classes, not just chosen representatives.
Recall that replacing $a$ by $a+kn$ changes neither the sum class nor
the product class. This is the well-definedness behind modular arithmetic.
Ask for the two reductions on the slide before displaying the next page.
}

\teachslide{The ring rules}{3}{
A \key{commutative ring with identity} has addition and multiplication.
\begin{itemize}
\item Addition is associative and commutative; it has $0$ and negatives.
\item Multiplication is associative and commutative; it has $1\ne0$.
\item Multiplication distributes over addition.
\end{itemize}
\[
a(b+c)=ab+ac.
\]
}{
State closure under both operations. Say explicitly that $a+0=a$,
$a+(-a)=0$, and $a\cdot1=a$. Associativity allows us to omit parentheses.
Integers and residue classes are the running examples; real numbers also
form a ring. These are the commutative rings with identity used in the
course; other conventions exist. Do not add the field inverse axiom yet.
The calculations inherited from integer arithmetic justify the axioms
for $\Z/n\Z$.
}

\teachslide{A new phenomenon: zero divisors}{2}{
In $\Z/6\Z$,
\[
\key{2\cdot3=0,\qquad 2\ne0,\quad3\ne0.}
\]
A \key{zero divisor} is a nonzero element $a$ for which
\[
ab=0\quad\text{for some nonzero }b.
\]
Cancellation can fail:
\[
2\cdot1=2\cdot4,\qquad 1\ne4.
\]
}{
Check both products in the last equation: $2$ and $8$ have the same
residue modulo $6$. Link cancellation failure to $2(1-4)=0$.
Ask whether this could happen in the rational numbers. The point is
that the arithmetic laws remain valid while familiar cancellation
requires extra hypotheses.
}

\teachslide{Which elements can we divide by?}{2}{
A \key{unit} is an element $u$ with a multiplicative inverse:
\[
uv=1\quad\text{for some }v.
\]
In $\Z/n\Z$,
\[
\key{[a]\text{ is a unit}\iff\gcd(a,n)=1.}
\]
In $\Z/6\Z$, the units are $1$ and $5$.
}{
Recall Bezout: $ax+ny=1$ supplies the inverse $[x]$. Conversely,
$ax\equiv1\pmod n$ gives an integer equation $ax+ny=1$, forcing the gcd
to be $1$. Thus inverses are not a new computational trick: extended
Euclid already computes them. A unit cannot be a zero divisor, since
multiplication by its inverse would force $b=0$.
}

\teachslide{Your turn: modulo eight}{3}{
Work in $\Z/8\Z$.
\begin{enumerate}
\item Which of $1,2,\ldots,7$ are units?
\item Find the inverse of $3$.
\item Find two nonzero elements whose product is $0$.
\end{enumerate}
\spacer
\key{Take one minute; then compare answers.}
}{
Keep this page visible while students calculate. Give a minute of
quiet work, then invite answers in chat or aloud. If needed, suggest
using the gcd test rather than trying every product. Ask students to
reduce $3^2$ and $2\cdot4$. The following page contains the answers.
}

\teachslide{Modulo eight: answers}{1}{
\[
\text{Units: }1,3,5,7.
\]
\[
3^{-1}=3,\qquad 3\cdot3=9=1\pmod8.
\]
\[
2\cdot4=0\pmod8.
\]
\key{Being nonzero does not guarantee an inverse.}
}{
Confirm that every odd residue is coprime to $8$. The nonzero even
residues are zero divisors. This is also a useful reminder that an
element may be its own inverse without equaling $1$.
}

\teachslide{A field allows division}{2}{
A \key{field} is a commutative ring with identity in which
\[
\text{every nonzero element has a multiplicative inverse.}
\]
Examples: $\Q$, $\R$, and $\Z/7\Z$.

A field has no zero divisors:
\[
ab=0,\ a\ne0\quad\Longrightarrow\quad b=a^{-1}0=0.
\]
}{
Division by $a\ne0$ means multiplication by $a^{-1}$. Zero has no
inverse because $0b=0\ne1$. Integers are a ring but are not a field;
for example $2$ has no integer inverse. Do not claim the converse
that every ring without zero divisors is a field: the integers refute it.
}

\teachslide{Why prime moduli give fields}{3}{
\[
\key{\Z/n\Z\text{ is a field}\iff n\text{ is prime}.}
\]
If $n=p$ is prime, then $1\le a<p$ implies $\gcd(a,p)=1$.

So every nonzero class has an inverse.
\spacer
If $n=rs$ with $1<r,s<n$, then
\[
[r][s]=[0]\qquad\text{with }[r],[s]\ne[0].
\]
We write $\F_p=\Z/p\Z$.
}{
Give both directions. In the composite case, $r$ cannot be invertible,
since an inverse would turn $rs=0$ into $s=0$. The statement assumes
$n\ge2$. Use $p=7$ and $n=8$ as the examples students have just seen.
Introduce the blackboard-bold F notation slowly: this is the same
familiar residue arithmetic, now recognized as a field.
}

\teachslide{Characteristic: count copies of one}{2}{
Write $m\cdot1$ for the sum of $m$ copies of $1$.

The \key{characteristic} is the least positive $m$ such that
\[
m\cdot1=0.
\]
If no such $m$ exists, the characteristic is $0$.
\[
\operatorname{char}(\Q)=0,\qquad
\operatorname{char}(\Z/n\Z)=n.
\]
}{
Distinguish repeated addition from an external scalar multiplication
that has not yet been defined. Characteristic zero means that no
positive number of copies of $1$ is zero, not that zero copies sum to
zero. In residues the least positive return is exactly the modulus.
Characteristic is defined for our rings as well as fields.
}

\teachslide{A field's positive characteristic is prime}{3}{
Suppose the characteristic of a field is $m=rs$, with $1<r,s<m$.
\[
(r\cdot1)(s\cdot1)=m\cdot1=0.
\]
Minimality of $m$ gives
\[
r\cdot1\ne0,\qquad s\cdot1\ne0.
\]
This contradicts the absence of zero divisors.
\spacer
\key{A field has characteristic $0$ or a prime.}
}{
Use distributivity to explain why the product of the two repeated sums
is the repeated sum with $rs$ terms. The minimality step is essential:
neither factor is zero since both indices are smaller than $m$.
Remind students that $m=1$ is excluded by $1\ne0$. The ring
$\Z/6\Z$ shows why the field hypothesis is needed.
}

\teachslide{A finite field cannot have characteristic zero}{2}{
In a finite field $K$, the list
\[
0,\ 1,\ 2\cdot1,\ 3\cdot1,\ldots
\]
must contain a repetition.

If $r\cdot1=s\cdot1$ with $r<s$, subtract:
\[
(s-r)\cdot1=0.
\]
\key{Therefore $\operatorname{char}(K)=p$ for some prime $p$.}
}{
This uses finiteness and subtraction, then the preceding theorem.
Do not assume the field's number of elements equals its characteristic.
That tempting assumption is exactly what Lecture 2 will disprove.
Here we know only that the characteristic is some positive prime.
}

\teachslide{Every finite field contains a prime field}{2}{
If $\operatorname{char}(K)=p$, the elements
\[
0,\ 1,\ 2\cdot1,\ldots,(p-1)\cdot1
\]
form a copy of $\F_p$ inside $K$.

Addition and multiplication agree with arithmetic modulo $p$.
\spacer
\key{$K$ may have more elements than these $p$ elements.}
}{
The map $[a]\mapsto a\cdot1$ is well defined modulo $p$, respects the
two operations, and is injective by minimality of the characteristic.
The nonzero elements have inverses within this copy by the same Bezout
argument as before. This copy is called the prime subfield. Prepare
the next step: elements of this smaller field can act as scalars on $K$.
}

\teachslide{Vectors and scalars}{2}{
A \key{vector space over a field $k$} has
\begin{itemize}
\item vectors $u,v,\ldots$ that we can add;
\item scalars $a,b,\ldots\in k$ that multiply vectors.
\end{itemize}
\[
v\longmapsto av.
\]
Examples: $\R^2$ over $\R$ and $\F_3^2$ over $\F_3$.
}{
Always name both the vector space and its scalar field. Vectors need
not be arrows and scalars need not be real numbers. For now use pairs
to keep every operation visible. Later the vectors will be elements
of a larger field. Scalar multiplication is an operation with one
scalar input and one vector input, rather than two vector inputs.
}

\teachslide{The vector-space rules}{3}{
Vector addition has the usual associative and commutative laws,
a zero vector, and additive inverses.

For $a,b\in k$ and $u,v\in V$,
\[
\begin{aligned}
a(u+v)&=au+av, & (a+b)v&=av+bv,\\
a(bv)&=(ab)v, & 1v&=v.
\end{aligned}
\]
\key{The field $k$ supplies the scalar arithmetic.}
}{
State closure under vector addition and scalar multiplication.
Read each law aloud and say which symbols are scalars and which are
vectors. The zero vector is distinct in type from the scalar $0$,
even when we use the same symbol. In coordinate spaces these axioms
hold because the field laws hold in every coordinate. This is a
definition, not an assertion that vectors can be multiplied together.
}

\teachslide{A familiar calculation over a different field}{2}{
In $\F_3^2$, let $u=(1,2)$ and $v=(2,2)$.
\[
u+v=(0,1),\qquad 2u=(2,1).
\]
\[
2u+v=(1,0).
\]
All coordinates are reduced modulo $3$.
\spacer
\key{Same vector rules; different scalar arithmetic.}
}{
Calculate each coordinate before reducing: $(3,4)$, $(2,4)$, and
$(4,3)$. Ask students to verify $2(u+v)=2u+2v=(0,2)$.
Three copies of any vector in this example sum to the zero vector.
Contrast that with $\R^2$ without relying on a geometric picture.
}

\teachslide{Linear combinations and span}{3}{
A \key{linear combination} of $v_1,\ldots,v_r$ is
\[
a_1v_1+\cdots+a_rv_r,\qquad a_i\in k.
\]
Their \key{span} is the set of all such combinations.

In $\F_3^2$, for $u=(1,2)$,
\[
\Span(u)=\{(0,0),(1,2),(2,1)\}.
\]
}{
There are three scalar choices here. The span includes zero and the
original vector, and is closed under addition and scalar multiplication.
For pairs $e_1=(1,0)$, $e_2=(0,1)$, every $(a,b)$ equals $ae_1+be_2$.
Mention this as a contrast: one vector need not span the whole space.
Independence and uniqueness of coordinates are Lecture 2's task.
}

\teachslide{Your turn: a span over $\F_2$}{3}{
In $\F_2^2$, let $u=(1,1)$ and $v=(1,0)$.
\begin{enumerate}
\item List the vectors in $\Span(u)$.
\item List all vectors $au+bv$ with $a,b\in\F_2$.
\item Can $(0,1)$ be written as $au+bv$?
\end{enumerate}
\spacer
\key{Remember: the only scalars are $0$ and $1$.}
}{
Give at least a minute of quiet calculation. Encourage a four-row
table for $(a,b)=(0,0),(1,0),(0,1),(1,1)$. For the third question,
ask for coefficients, not merely a yes. This exercise prepares the
coordinate-counting argument without introducing independence early.
}

\teachslide{Four coefficient choices, four vectors}{1}{
\[
\begin{array}{c|c}
(a,b)&au+bv\\\hline
(0,0)&(0,0)\\
(1,0)&(1,1)\\
(0,1)&(1,0)\\
(1,1)&(0,1)
\end{array}
\]
\[
\Span(u)=\{(0,0),(1,1)\},\qquad (0,1)=u+v.
\]
}{
The table shows both spanning and uniqueness in this example.
We will name the combination of those properties a basis next time.
Do not generalize distinct coefficient choices giving distinct
vectors to an arbitrary spanning list: independence is needed.
}

\teachslide{What we have built today}{2}{
\begin{itemize}
\item Residue arithmetic gives rings; prime moduli give fields.
\item Every finite field contains some $\F_p$.
\item Vector spaces can use $\F_p$ as their scalar field.
\end{itemize}
\spacer
Next time: \key{bases, counting, and a field with $49$ elements.}
}{
Ask students to explain the distinction between a unit and a nonzero
element. Then connect the three bullets: the copy of $\F_p$ supplies
scalars for a vector space whose vectors are elements of the larger
field. We will check its vector-space rules next time rather than
asking them to accept the slogan on faith.
}

\teachslide{Exit question}{2}{
Why is $\Z/6\Z$ a ring but not a field?

Why can it not serve as the scalar \emph{field} in our definition
of a vector space?
\spacer
\key{Give a specific arithmetic obstruction.}
}{
Expected answer: $2\cdot3=0$ with both factors nonzero, so $2$ and $3$
cannot have inverses. Our vector-space definition requires a field
of scalars. Addition and scalar multiplication over a ring can still
be studied; those objects are modules, an optional observation on the
backup slide. End here at minute 50.
}

\backups
\backupslide{Rings as scalars: a glimpse beyond vector spaces}{
One can use a ring of scalars and keep the vector-space identities.

The resulting object is called a \key{module}.

Over $\Z/6\Z$, $2v=0$ need not imply $v=0$:
\[
v=3,\qquad 2v=0.
\]
Our next lectures use fields, where nonzero scalars are invertible.
}{
Use only if a student asks whether the ring case is meaningless.
It is meaningful, but some basis and dimension facts change. The
example uses the ring itself as a module over itself. Avoid developing
module theory during the main lecture.
}

\backupslide{One set, two different scalar fields}{
The complex numbers form a vector space over $\R$:
\[
z=a+bi,\qquad a,b\in\R.
\]
The same set forms a vector space over $\mathbb C$:
\[
z=z\cdot1.
\]
\key{Changing the scalars can change the dimension.}
}{
This example is optional because complex numbers may not be familiar.
Over the reals, $1,i$ form a basis; over the complex numbers, $1$ alone
does. Lecture 2 supplies the precise basis definition. The example
anticipates $\F_{49}$ over $\F_7$ without implying any embedding into
the ordinary real or complex fields.
}

\clearpage
\setcounter{slide}{0}\setcounter{elapsed}{0}
\section*{Lecture 2: Counting and building finite fields}
\teachcover{Counting and building\\finite fields}{1}{
Bases turn algebra into counting
\[
\F_7\subset\F_{49}
\]
}{
This lecture finishes the vector-space reminder, proves the necessary
prime-power size theorem, and constructs the quadratic extension needed
for Fibonacci modulo $7$. The construction proves existence for this
example, not for every prime power. The main route now proves that
dimension is well defined by a six-slide exchange argument. Keep the
basis calculation break; the extra extension-arithmetic exercise is
optional material after the exit question.
}

\teachslide{Recall: scalars determine the combinations}{1}{
In $\F_3^2$, let $u=(1,2)$.
\[
\Span(u)=\{0u,1u,2u\}
=\{(0,0),(1,2),(2,1)\}.
\]
With $e_1=(1,0)$ and $e_2=(0,1)$,
\[
(a,b)=ae_1+be_2.
\]
\key{How can we tell whether coordinates are unique?}
}{
Retrieve rather than re-lecture: ask for $2u$, and ask how to express
$(2,1)$ in the standard coordinate vectors. A spanning list expresses
every vector, but could express the same vector in different ways.
The next two slides isolate that issue.
}

\teachslide{A redundant spanning list}{2}{
In $\F_3^2$, put $w=e_1+e_2$.

The list $e_1,e_2,w$ spans the space, but
\[
w=0e_1+0e_2+1w=1e_1+1e_2+0w.
\]
The difference gives a nontrivial relation:
\[
\key{e_1+e_2-w=0.}
\]
}{
Point out two distinct coefficient triples for the same vector.
Modulo $3$, the coefficients in the relation are $1,1,2$, so they
are not all zero. The problem is redundancy, not lack of spanning.
Ask which vector can be removed without changing the span.
}

\teachslide{Linear independence}{3}{
Vectors $v_1,\ldots,v_d$ are \key{linearly independent} if
\[
a_1v_1+\cdots+a_dv_d=0
\quad\Longrightarrow\quad a_1=\cdots=a_d=0.
\]
Then two expressions for the same vector have the same coefficients:
\[
\sum a_iv_i=\sum b_iv_i
\ \Longrightarrow\ \sum(a_i-b_i)v_i=0.
\]
}{
Finish the second line verbally: independence forces every $a_i-b_i=0$.
Conversely, unique expressions of zero force independence. Give the
standard pair $e_1,e_2$ as the basic example: $(a,b)=0$ implies both
coordinates are zero. Distinguish a nontrivial relation from one whose
coefficients all vanish.
}

\teachslide{A basis gives existence and uniqueness}{2}{
A \key{basis} is a linearly independent spanning list.

Every $v\in V$ has a unique expression
\[
\key{v=a_1v_1+\cdots+a_dv_d.}
\]
\begin{itemize}
\item Spanning gives existence of the coefficients.
\item Independence gives uniqueness.
\end{itemize}
}{
Have students supply which hypothesis does which job. An ordered basis
also specifies the order of coordinates. The standard basis of $k^d$
is the list of coordinate vectors. We will count ordered bases when
we count matrices in Lecture 3.
}

\teachslide{Build a basis one vector at a time}{3}{
Start with no vectors; their span is $\{0\}$.

If the current span is not all of $V$, choose a vector outside it.

Adding that vector preserves independence.
\spacer
\key{In a finite space, this process must stop.}

At stopping, the chosen vectors form a basis.
}{
Prove the key step. If $a_1v_1+\cdots+a_rv_r+bw=0$ and $b\ne0$,
then division by $b$ expresses $w$ in the old span, contradicting its
choice. Thus $b=0$, and old independence makes all other coefficients
zero. Finiteness makes the strictly increasing sequence of spans stop.
This proof specifically establishes the finite basis needed below.
}

\teachslide{Could two bases have different lengths?}{1}{
A basis gives unique coordinates.

But before defining dimension, we must prove that its length
does not depend on the chosen basis.
\spacer
Our key comparison:
\[
\key{\text{independent list length}\ \le\ \text{spanning list length}.}
\]
}{
Make the logical issue explicit. Existence and uniqueness of
coordinates refer to one chosen basis. They do not yet compare two
different bases. The following argument works over any field,
including infinite fields; it does not use the count $p^d$.
Only finite lists are involved in the comparison theorem.
}

\teachslide{Exchange one vector}{2}{
Suppose $w_1,\ldots,w_m$ span $V$ and $v\ne0$.
\[
v=a_1w_1+\cdots+a_mw_m.
\]
Some coefficient is nonzero; relabel so $a_1\ne0$.
\[
\key{w_1=a_1^{-1}\bigl(v-a_2w_2-\cdots-a_mw_m\bigr).}
\]
Thus $v,w_2,\ldots,w_m$ still span $V$.
}{
Explain both inclusions of spans. Each vector in the new list lies
in the old span. The solved equation puts the discarded $w_1$ in
the new span, and all other old vectors remain. Hence the spans are
equal. The spanning list need not be independent. Relabeling just
moves a nonzero coefficient to the first position. Dividing by it
is valid because the scalars form a field.
}

\teachslide{An exchange in $\F_3^2$}{2}{
Start with $e_1=(1,0)$ and $e_2=(0,1)$.

Insert $u=(1,1)=e_1+e_2$.
\[
e_1=u-e_2,\qquad \key{u,e_2\text{ still span}.}
\]
Now insert $v=(1,2)=u+e_2$.
\[
e_2=v-u,\qquad \key{u,v\text{ still span}.}
\]
}{
Ask students to supply each subtraction. Over $\F_3$ a minus sign
means adding twice that vector. Each exchange preserves the number
of vectors and the span. The general proof will also need to ensure
we can always remove an old $w$, rather than a vector already inserted
from the independent list. This example anticipates the retained
basis exercise later in the lecture.
}

\teachslide{Repeat the exchange}{3}{
Let $v_1,\ldots,v_r$ be independent and $w_1,\ldots,w_m$ span $V$.

After inserting $v_1,\ldots,v_k$, the list still spans and has
\[
v_1,\ldots,v_k,\quad m-k\text{ remaining old vectors}.
\]
Express $v_{k+1}$ in this list.

\key{Some remaining old vector has a nonzero coefficient.}

Otherwise $v_{k+1}\in\Span(v_1,\ldots,v_k)$, contradicting independence.
}{
Start with $k=0$. Since the list spans, such an expression exists.
If all coefficients of the remaining old vectors vanished, moving
$v_{k+1}$ to the other side would give a relation among the $v_i$
with its coefficient -1, contradicting independence. Therefore
solve for an old vector having nonzero coefficient and replace that
vector by $v_{k+1}$. The earlier exchange formula preserves spanning
and never removes an already inserted $v_i$.
}

\teachslide{An independent list cannot be longer}{2}{
Every insertion replaces one remaining old vector.

After $m$ insertions, the spanning list would be
\[
v_1,\ldots,v_m.
\]
If $r>m$, then $v_{m+1}$ would be in their span, contradicting independence.
\[
\key{r\le m.}
\]
This is the \key{Steinitz exchange argument}.
}{
The contradiction is the rigorous reason the process cannot outlast
the original spanning list. If fewer than $m$ independent vectors
were given, the process stops earlier and the inequality already
holds. If $m=0$, the space is $\{0\}$ and any nonempty independent
list is impossible, so the assertion remains true. The argument
also shows an independent finite list can be extended to a basis
when a finite spanning list is available: keep the inserted vectors
and remove any remaining redundancies.
}

\teachslide{Dimension is well defined}{2}{
Suppose $B$ and $C$ are bases, of lengths $r$ and $m$.
\begin{itemize}
\item $B$ independent, $C$ spanning: $r\le m$.
\item $C$ independent, $B$ spanning: $m\le r$.
\end{itemize}
Therefore \key{$r=m$}.

For a space with a finite basis, this common length is its dimension:
\[
\dim_k V=r.
\]
}{
Only now define dimension. Such a space is called finite-dimensional.
The result proves independence from the choice of basis, not
independence from the choice of scalar field. The theorem compares
finite bases. It also rules out an infinite independent basis when
a finite spanning list exists, since an arbitrarily long finite
sublist would violate the inequality. No Gaussian elimination was
needed, though elimination supplies another proof of the same
comparison theorem. The zero space has the empty basis and dimension
zero. The next slide specializes to finite-field coordinates.
}


\teachslide{Count coordinate choices}{3}{
Suppose $v_1,\ldots,v_d$ is a basis over $\F_p$.
\[
v=a_1v_1+\cdots+a_dv_d,
\qquad a_i\in\F_p.
\]
Each of the $d$ coordinates has $p$ choices.

Uniqueness means that different choices give different vectors.
\[
\key{|V|=p^d.}
\]
}{
The coordinate map from $\F_p^d$ to $V$ is onto by spanning and
one-to-one by independence. Thus the product rule counts the vectors
exactly, not merely bounds their number. In particular, a span of
$r$ independent vectors has $p^r$ elements. If two bases have lengths
$d,e$, this count gives $p^d=p^e$ and hence $d=e$.
}

\teachslide{Your turn: is this a basis?}{3}{
In $\F_3^2$, let $u=(1,1)$ and $v=(1,2)$.
\begin{enumerate}
\item Solve $au+bv=(0,0)$.
\item Express $(2,0)$ as $au+bv$.
\item How many vectors are in $\Span(u,v)$?
\end{enumerate}
\spacer
\key{Write the two coordinate equations.}
}{
Give time to write $a+b=0$ and $a+2b=0$. Subtract the first from
the second to get $b=0$, then $a=0$. For $(2,0)$, use $a+b=2$,
$a+2b=0$. A suggestion for students stuck in real-number habits:
$-2=1$ in $\F_3$. The answer appears on the next page.
}

\teachslide{Basis exercise: answers}{1}{
\[
a+b=0,\quad a+2b=0\quad\Longrightarrow\quad b=0, a=0.
\]
So $u,v$ are independent.
\[
\key{(2,0)=u+v.}
\]
Their span has $3^2=9$ elements, all of $\F_3^2$.

Thus $u,v$ form a basis.
}{
There are nine elements in the ambient space, so the nine-element
span is the whole space. We have checked both requirements for a
basis without silently assuming that independence implies spanning
in arbitrary spaces. Coordinates depend on the chosen basis:
$(2,0)$ has coordinates $(1,1)$ in this basis.
}

\teachslide{A field can be a vector space}{3}{
Let $K$ be a finite field of characteristic $p$.

It contains $\F_p$. Regard
\begin{itemize}
\item elements of $K$ as vectors;
\item elements of $\F_p$ as scalars;
\item multiplication in $K$ as scalar multiplication.
\end{itemize}
\[
a(u+v)=au+av,\qquad a(bu)=(ab)u.
\]
}{
Check all the vector-space axioms, not just the two displayed ones:
addition in $K$ is an abelian group operation, $(a+b)u=au+bu$, and
$1u=u$. These are field laws restricted to the scalar subfield.
This is not making $K$ a vector space over the real numbers.
Here a vector can also be multiplied by another vector because $K$
has additional field structure; the vector-space definition alone
does not supply that multiplication.
}

\teachslide{Every finite field has prime-power size}{2}{
Finiteness gives a basis of $K$ over its prime field $\F_p$.
\[
\key{|K|=p^d,\qquad p\text{ prime},\quad d\ge1.}
\]
\begin{itemize}
\item Characteristic gives the prime $p$.
\item A basis over $\F_p$ gives the exponent $d$.
\end{itemize}
\spacer
There is no field with $6$ elements.

A field with $49$ elements would have characteristic $7$.
}{
Ask students to rule out $6$ as a prime power. For $49$, unique prime
factorization gives $p=7,d=2$. We have proved a necessary condition;
we have not proved existence or uniqueness for arbitrary $p^d$.
The next construction supplies a particular missing example.
To justify a finite basis, apply the earlier basis-building process
to the finite set $K$. The strictly increasing spans must stop.
Since $1\ne0$, its dimension over $\F_p$ is at least one. Thus
finiteness first gave prime characteristic and now gives a finite basis.
}

\teachslide{Build the missing square root}{2}{
Over $\F_7$, $5$ is not a square.

Introduce a symbol $s$ with $s^2=5$, and use formal pairs
\[
\key{K=\{a+bs:a,b\in\F_7\}.}
\]
Two elements are equal exactly when their two coefficients agree.
\[
(a+bs)+(c+es)=(a+c)+(b+e)s.
\]
}{
Define the elements as pairs first, so uniqueness of this representation
is built in. The symbol $s$ is not an ordinary real square root.
We are defining new finite arithmetic. Write $s=(0,1)$ and embed
$a\in\F_7$ as $(a,0)$. There are already $7^2$ candidate elements;
we still must prove this arithmetic is a field.
}

\teachslide{Multiply, then replace $s^2$ by $5$}{3}{
\[
(a+bs)(c+es)=(ac+5be)+(ae+bc)s.
\]
For example,
\[
\begin{aligned}
(2+s)(3+2s)&=6+7s+2s^2\\
&=16+7s\\
&=2.
\end{aligned}
\]
All coefficients are reduced modulo $7$.
}{
Explain the four terms from distribution. This is polynomial arithmetic
with the relation $s^2=5$; the ring laws hold. If requested, associativity
can be checked from the pair formula, or justified by reducing products
of polynomials modulo $X^2-5$. A backup slide introduces that notation.
The product $2$ is an element of the embedded $\F_7$.
}

\teachslide{Why every nonzero element has an inverse}{3}{
\[
(a+bs)(a-bs)=a^2-5b^2.
\]
For $(a,b)\ne(0,0)$, this denominator is nonzero:
\begin{itemize}
\item If $b=0$, then $a^2\ne0$.
\item If $b\ne0$, its vanishing would give $(a/b)^2=5$.
\end{itemize}
Therefore
\[
\key{(a+bs)^{-1}=\frac{a-bs}{a^2-5b^2}.}
\]
}{
The second case contradicts that $5$ is nonsquare in $\F_7$. The
denominator is a nonzero element of the original field, so its inverse
exists there. This proves the new ring is a field. Explain the division
notation as multiplication by that inverse. If $5$ were a square, the
same formal construction could have nonzero zero divisors; a backup
slide gives the example modulo $11$.
}

\teachslide{A field with forty-nine elements}{2}{
We have built a field $K$, which we call $\F_{49}$.
\[
\key{\{1,s\}\text{ is a basis over }\F_7.}
\]
Every element is a unique $a\cdot1+b\cdot s$.
\[
|K|=7^2=49,\qquad \dim_{\F_7}K=2.
\]
Its characteristic is $7$.
}{
Spanning follows from the definition as pairs, and independence follows
because $a+bs=0$ means both coordinates are zero. Seven copies of $1$
sum to zero; no smaller positive number does. Thus the characteristic
is $7$, although the number of elements is $49$. The same construction
works over an odd prime with any nonsquare $d$ replacing $5$.
}

\teachslide{$\F_{49}$ and $\Z/49\Z$ are different}{2}{
\[
\begin{array}{c|c|c}
&\F_{49}&\Z/49\Z\\\hline
\text{Elements}&49&49\\
\text{Characteristic}&7&49\\
\text{Every nonzero element a unit?}&\text{Yes}&\text{No}
\end{array}
\]
In $\Z/49\Z$, the nonzero residue $7$ satisfies $7^2=0$.

In $\F_{49}$, seven copies of $1$ already give $0$.
}{
Read the two columns slowly. The symbol $7$ denotes different things
in the two arithmetics: a nonzero residue in one and the sum $7\cdot1=0$
in the other. Equal cardinality is far weaker than equal arithmetic.
This is the main misconception to address before moving on.
}

\teachslide{Exit question}{2}{
Suppose $K$ is a field with $25$ elements.
\begin{enumerate}
\item What is its characteristic?
\item What is its dimension over its prime field?
\item Why is $K$ not $\Z/25\Z$?
\end{enumerate}
}{
Expected answers: characteristic $5$, dimension $2$ over $\F_5$, and
$\Z/25\Z$ has characteristic $25$ and the nonzero zero divisor $5$.
The question assumes a field with $25$ elements; it does not rely on
a general finite-field existence theorem. This is the 50-minute stop.
}

\backups
\backupslide{Polynomial notation for the construction}{
\[
\key{\F_{49}=\F_7[X]/(X^2-5).}
\]
Two polynomials represent the same element when their difference
is a multiple of $X^2-5$.

Polynomial division gives a unique remainder $a+bX$.

The class of $X$ is our element $s$, and $s^2=5$.
}{
This is the polynomial analogue of integer residues. Addition and
multiplication descend from polynomial arithmetic, so the ring laws
follow. Polynomial division by the monic quadratic gives unique
remainders of degree at most one. The main lecture proves the field
property directly, so a general quotient-field theorem is unnecessary.
}

\backupslide{What if the chosen number is already a square?}{
Try the same formal-pair construction over $\F_{11}$, with $s^2=5$.

Since $4^2=5$,
\[
(s-4)(s+4)=s^2-16=0.
\]
Both factors are nonzero formal pairs.
\spacer
\key{This construction gives a ring with zero divisors, not a field.}
}{
The coefficient of $s$ is $1$ in both factors, so neither is the
zero pair. This makes the nonsquare hypothesis in the inverse proof
concrete. There is no need to construct an extension when the roots
already lie in the original field.
}

\backupslide{Your turn: calculate in the new arithmetic}{
In $K$, where $s^2=5$ and coefficients lie in $\F_7$, compute
\[
(1+2s)(3+s).
\]
Then compute
\[
(2+s)(2-s).
\]
\key{Distribute first; reduce afterward.}
}{
Allow a minute of quiet work. Ask students to show intermediate
expressions so that an error in reduction can be distinguished from
an error in expansion. The second calculation prepares the conjugate
inverse formula and should not be described as complex conjugation.
}


\backupslide{Two products}{
\[
(1+2s)(3+s)=3+7s+2s^2=13=6.
\]
\[
\key{(2+s)(2-s)=4-s^2=4-5=6.}
\]
Since $6^{-1}=6$ in $\F_7$,
\[
(2+s)^{-1}=6(2-s)=5+s.
\]
}{
Verify $6\cdot6=1\pmod7$. Multiplying the second identity by $6$
gives the inverse of this one element. Do not yet claim that every
element is invertible; the next slide proves that general statement.
For a quick check, $(2+s)(5+s)=15+7s=1$.
}


\backupslide{The roots now exist}{
In $\F_{49}$, division by $2$ is allowed, and
\[
\alpha=\frac{1+s}{2},\qquad\beta=\frac{1-s}{2}.
\]
\[
\alpha+\beta=1,\qquad\alpha\beta=-1.
\]
Both solve $X^2-X-1=0$.
\spacer
Next: \key{multiplicative orders, matrices, and the period.}
}{
Modulo $7$, $2^{-1}=4$, so $\alpha=4+4s$ and $\beta=4-4s$.
Their sum and product verify the factorization
$(X-\alpha)(X-\beta)=X^2-X-1$. The roots are distinct because $s\ne0$.
This reconnects the construction to the original motivating question.
Lecture 3 supplies cyclicity and matrix counts; Lecture 4 uses the
roots' powers to finish the period argument.
}


\backupslide{A bridge to matrices}{
An invertible matrix has columns that form an ordered basis.
\[
\key{\text{Counting invertible matrices = counting ordered bases.}}
\]
We will count the column choices one at a time.
}{
Ask students to remember how many elements are in the span of one
nonzero vector over $\F_p$: exactly $p$. That is the key number for
the next lecture's second column. This optional bridge previews
Lecture 3; the main route already ended at the exit question.
}

\backupslide{Gaussian elimination gives another proof}{
Suppose $w_1,\ldots,w_m$ span and $v_1,\ldots,v_r$ are independent.

Choose coefficients $v_i=\sum_{j=1}^m c_{ji}w_j$; let $C=(c_{ji})$.

If $r>m$, elimination in the homogeneous system $Cx=0$
leaves a free variable, hence a nonzero solution.
\[
\sum_{i=1}^r x_iv_i=\sum_{j=1}^m(Cx)_jw_j=0,
\]
contradicting independence.
}{
$C$ has $m$ rows and $r$ columns. Elimination produces at most $m$
pivots, so when $r>m$ a variable is free. Set one free variable to 1,
the others to 0, and solve backwards for the pivot variables. Row
operations preserve the solutions, and division by a nonzero pivot
is allowed over a field. The coefficients $c_{ji}$ need not be unique:
the original $w_j$ merely span, and any chosen expressions suffice.
Thus this proof does not assume the dimension theorem or rank-nullity.
It is an alternative to the main exchange proof, not a prerequisite.
}


\clearpage
\setcounter{slide}{0}\setcounter{elapsed}{0}
\section*{Lecture 3: Cyclic groups and counting matrices}
\teachcover{Cyclic groups\\and counting matrices}{1}{
Multiplication in a finite field has a generator
\[
K^\times=\{1,g,g^2,\ldots,g^{q-2}\},\qquad |K|=q.
\]
}{
This lecture proves cyclicity of a finite field's multiplicative group
and then counts ordered bases and determinant fibers. Recall Lagrange's
theorem and the definition of element order. We prove the polynomial
root bound before using it. The group-size-one case is included.
The optional slides compare finite fields with units of prime-power
residue rings; they are not part of the 50-minute main route.
}

\teachslide{The nonzero elements form a group}{2}{
Let $K$ be a field with $q$ elements.
\[
K^\times=K\setminus\{0\},\qquad |K^\times|=N=q-1.
\]
Multiplication is commutative; the identity is $1$;
every element has an inverse.
\[
\key{\ord(a)\mid N,\qquad a^N=1\quad(a\in K^\times).}
\]
}{
Closure follows from absence of zero divisors: the product of two
nonzero elements cannot be zero. Inverses are nonzero. Lagrange
gives the displayed order divisibility. Define order as the least
positive exponent returning to $1$. The task is to find one element
of order $N$, which is stronger than all orders dividing $N$.
}

\teachslide{What cyclic means}{2}{
In $\F_5$, the powers of $2$ are
\[
\begin{array}{c|rrrr}
t&0&1&2&3\\\hline
2^t&1&2&4&3
\end{array}
\qquad 2^4=1.
\]
Every nonzero element is a power of $2$.
\spacer
A group is \key{cyclic} if one element generates the whole group.
}{
Calculate the last power by reducing $8$ modulo $5$. A generator of
a finite group of size $N$ has order $N$, and conversely an element
of order $N$ supplies all $N$ elements by its distinct powers.
There may be several generators. Our aim is to prove at least one
exists in every finite field, not to identify a canonical generator.
}

\teachslide{A polynomial cannot have too many roots}{2}{
Over a field, if $f(a)=0$, then
\[
f(X)=(X-a)h(X).
\]
Any other root $b\ne a$ is a root of $h$:
\[
0=(b-a)h(b)\quad\Longrightarrow\quad h(b)=0.
\]
Induction gives:
\[
\key{\deg f=d,\ f\ne0\quad\Longrightarrow\quad\text{at most }d\text{ roots}.}
\]
}{
Recall polynomial division by $X-a$. The constant nonzero case has
no roots. Each root factors off one linear term; cancellation works
since $b-a$ is a nonzero field element. This proof is the specific
place field arithmetic enters the cyclicity theorem. The root bound
can fail over a ring with zero divisors.
}

\teachslide{Build the order one prime at a time}{2}{
Factor $N=q-1$:
\[
N=\ell_1^{r_1}\cdots\ell_t^{r_t}.
\]
We will construct, for each prime factor $\ell$,
\[
\key{g_\ell\text{ with order }\ell^r,\qquad \ell^r\parallel N.}
\]
Then multiply these elements.

If $N=1$, the group $\{1\}$ is already cyclic.
}{
Explain the double-bar notation: $\ell^r$ divides $N$, but the next
power does not. Distinct primes have relatively prime powers. The
product step requires both coprime orders and commutativity, proved
below. For $K=\F_2$, $N=1$ and $1$ is a generator of the trivial group.
}

\teachslide{Find an element with the required prime factor}{3}{
Fix $\ell^r\parallel N$.

The polynomial
\[
X^{N/\ell}-1
\]
has at most $N/\ell<N$ roots.

It cannot vanish at every element of $K^\times$.
\[
\key{\text{Choose }a\in K^\times\text{ with }a^{N/\ell}\ne1.}
\]
}{
Do not assume a generator to choose $a$: the root bound proves that
at least one of the $N$ nonzero elements fails this equation. The
chosen $a$ can depend on $\ell$, which is fine. Its order must include
the full $\ell$-power of $N$, but the next slide proves the needed
statement directly through powers rather than invoking an order formula.
}

\teachslide{Isolate the prime-power order}{3}{
Set $M=N/\ell^r$ and $g_\ell=a^M$.
\[
g_\ell^{\ell^r}=a^N=1,
\qquad g_\ell^{\ell^{r-1}}=a^{N/\ell}\ne1.
\]
Its order divides $\ell^r$, but does not divide $\ell^{r-1}$.

Every proper divisor of $\ell^r$ divides $\ell^{r-1}$.
\[
\key{\ord(g_\ell)=\ell^r.}
\]
}{
Read both exponent products aloud. If $x^m=1$, its order divides
$m$: divide $m$ by the least positive return time and use minimality
of the remainder. This fact is more precise than Lagrange and also
applies here to a known exponent. The prime-power divisor observation
is what forces the exact order without computing it by trial.
}

\teachslide{Combine two coprime orders}{4}{
Suppose $a,b$ commute, with orders $r,s$ and $\gcd(r,s)=1$.
\[
(ab)^{rs}=1.
\]
If $(ab)^t=1$, put $z=a^t=b^{-t}$. Then
\[
z^r=z^s=1\quad\Longrightarrow\quad z=1
\]
by Bezout's identity. Thus $r\mid t$ and $s\mid t$, so $rs\mid t$.
\[
\key{\ord(ab)=rs.}
\]
}{
This lemma deserves the full four minutes. If $ur+vs=1$, then
$z=z^{ur+vs}=(z^r)^u(z^s)^v=1$; negative exponents use inverses.
Hence $a^t=b^t=1$. The order-divides-return-time fact forces both
divisibilities, whose coprimality gives $rs\mid t$. We already know
$rs$ is a return time, so it is the least. Commutativity is needed
for $(ab)^t=a^tb^t$. The result does not claim that arbitrary
products have the lcm of their factors' orders.
}

\teachslide{One element of full order}{3}{
The prime-power orders $\ell_i^{r_i}$ are pairwise coprime.

By repeatedly applying the lemma,
\[
g=g_{\ell_1}\cdots g_{\ell_t}
\quad\text{has order}\quad
\ell_1^{r_1}\cdots\ell_t^{r_t}=N.
\]
Its $N$ distinct powers exhaust $K^\times$.
\[
\key{\text{The multiplicative group of every finite field is cyclic.}}
\]
}{
Check that a partial product's order remains coprime to the next
prime-power order, so induction applies. An element of order $N$
has exactly $N$ distinct powers from exponent $0$ through $N-1$;
there are only $N$ elements in the group. The proof has used Lagrange,
the polynomial root bound, Bezout, and commutativity, with no general
classification theorem for finite abelian groups.
}

\teachslide{Primitive elements}{2}{
A generator $g$ of $K^\times$ is called a \key{primitive element}.
\[
K^\times=\{1,g,\ldots,g^{q-2}\},\qquad \ord(g)=q-1.
\]
If $K$ has characteristic $p$, then
\[
\key{K=\F_p(g):\text{ the field is generated by one element}.}
\]
Every nonzero element is a power of $g$.
}{
In this finite-field context primitive means generator of the
multiplicative group. Such an element also generates the extension
field, since a field containing it contains all its powers and zero.
The converse need not hold: generating a field extension does not
necessarily mean having multiplicative order $q-1$. In particular,
the Fibonacci root modulo $7$ will have order $16$, not $48$.
}

\teachslide{Your turn: combine orders in $\F_7$}{3}{
In $\F_7^\times$,
\[
\ord(6)=2,\qquad\ord(4)=3.
\]
Verify these orders, then answer:
\begin{enumerate}
\item What is the order of $6\cdot4$?
\item Which residue is therefore a primitive element?
\end{enumerate}
}{
Give a minute of calculation. Students should find $6^2=1$ with
$6\ne1$, and $4^2=2$, $4^3=1$ with $4\ne1$. Coprimality lets us
combine the orders. This is a concrete version of the theorem's
final construction, not a search through all six possible generators.
}

\teachslide{A generator modulo seven}{1}{
\[
6\cdot4=3\pmod7,
\qquad \ord(3)=2\cdot3=6.
\]
\[
\begin{array}{c|rrrrrr}
t&0&1&2&3&4&5\\\hline
3^t&1&3&2&6&4&5
\end{array}
\]
\key{One element generates all six nonzero residues.}
}{
The table verifies the conclusion independently. We used the lemma
to know the order before listing all powers. This ends the cyclicity
portion at minute 28. The rest of this lecture counts matrix groups.
}


\teachslide{From field elements to matrices}{2}{
Over $\F_p$, the span of $r$ independent vectors has $p^r$ elements.
\spacer
An invertible matrix has columns forming an ordered basis.
\[
\key{\text{Count invertible matrices by choosing their columns.}}
\]
}{
Retrieve the coordinate count from Lecture 2. The remaining 20 minutes
count GL and SL, providing the group-size bounds we will use for the
Fibonacci matrix in Lecture 4. The cyclicity theorem concerns field
elements under multiplication, not a claim that GL or SL is cyclic.
}

\teachslide{Invertible matrices and ordered bases}{2}{
For a matrix with columns $v_1,\ldots,v_d$,
\[
M\begin{pmatrix}a_1\\\vdots\\a_d\end{pmatrix}
=a_1v_1+\cdots+a_dv_d.
\]
$M$ is invertible exactly when every output has unique coordinates.
\spacer
\key{Its columns must form an ordered basis of $\F_p^d$.}
}{
Connect matrix invertibility to the bijectivity of its linear map.
A basis supplies existence and uniqueness of solutions to $Mx=y$.
Conversely, an invertible map makes the columns spanning and independent.
Over a field this also agrees with nonzero determinant. Order matters:
exchanging columns usually gives another matrix.
}

\teachslide{Choose the first column}{2}{
In $\F_p^2$, there are $p^2$ vectors.

The first basis vector can be anything except zero.
\[
\key{p^2-1\text{ choices}.}
\]
\[
v_1\ne\begin{pmatrix}0\\0\end{pmatrix}.
\]
}{
Ask students why a zero column prevents invertibility. A singleton
list containing a nonzero vector is independent because a nonzero
scalar has an inverse. This will matter when we count its span.
}

\teachslide{Choose the second column}{3}{
Once $v_1\ne0$ is fixed, its span has $p$ elements:
\[
\Span(v_1)=\{av_1:a\in\F_p\}.
\]
Choose $v_2$ outside that span:
\[
\key{p^2-p\text{ choices}.}
\]
Therefore
\[
|\GL_2(\F_p)|=(p^2-1)(p^2-p).
\]
}{
The scalar multiples are distinct because $v_1$ is nonzero. Being
outside the span gives independence, by the basis-extension argument
from Lecture 2. Two independent vectors in $\F_p^2$ span it: their
span has $p^2$ elements, the full ambient size. The number of allowed
second columns does not depend on which first column was chosen.
}

\teachslide{Your turn: count modulo two}{2}{
How many invertible $2\times2$ matrices are there over $\F_2$?

Fix the first column to be
\[
v_1=\begin{pmatrix}1\\0\end{pmatrix}.
\]
Which second columns are allowed?
}{
Give quiet work time, then ask for both the total count and the actual
two allowed columns for this fixed choice. Students should list the
four vectors and exclude $0$ and $v_1$. This reinforces the reason
for subtraction in the count.
}

\teachslide{Six invertible matrices}{1}{
\[
|\GL_2(\F_2)|=(4-1)(4-2)=6.
\]
With $v_1=(1,0)^{\mathsf T}$, the allowed second columns are
\[
\begin{pmatrix}0\\1\end{pmatrix},\qquad
\begin{pmatrix}1\\1\end{pmatrix}.
\]
The span we excluded is $\{(0,0),(1,0)\}$.
}{
Confirm that zero and the first column are exactly the forbidden
vectors. Over $\F_2$, every nonzero determinant is $1$, so GL and
SL will have the same size. Mention this only as an anticipation.
}

\teachslide{The same count in dimension $d$}{2}{
After $r$ independent columns, the span has $p^r$ vectors.

The next column has $p^d-p^r$ choices.
\[
\key{|\GL_d(\F_p)|=\prod_{r=0}^{d-1}(p^d-p^r).}
\]
The first factor is $p^d-1$; the last is $p^d-p^{d-1}$.
}{
The span of the empty list has size $1$, accounting for the first
factor when $r=0$. Multiplying the independent choices counts ordered
bases. Do not expand or simplify the product: the unsimplified form
records the argument. Most lecture time stays with $d=2$.
}

\teachslide{The determinant groups matrices by value}{2}{
\[
\det:\GL_d(\F_p)\longrightarrow\F_p^\times.
\]
\[
\det(MN)=\det(M)\det(N).
\]
Every nonzero $a\in\F_p$ occurs:
\[
D_a=\operatorname{diag}(a,1,\ldots,1),\qquad\det D_a=a.
\]
\key{$\SL_d(\F_p)$ is the set with determinant $1$.}
}{
Define $\F_p^\times$ as the multiplicative group of nonzero field
elements, of size $p-1$. The determinant is a surjective group
homomorphism. Quotient-group notation is optional; the next slide
uses a concrete bijection between determinant fibers.
}

\teachslide{Every determinant fiber has the same size}{2}{
Multiplication by $D_a$ gives a bijection
\[
\SL_d(\F_p)\longrightarrow\{M:\det M=a\},
\qquad S\longmapsto D_aS.
\]
The inverse is $M\longmapsto D_a^{-1}M$.

There are $p-1$ nonzero determinant values, so
\[
\key{|\SL_d(\F_p)|=\frac{|\GL_d(\F_p)|}{p-1}.}
\]
}{
Verify the two determinants: $\det(D_aS)=a$ and
$\det(D_a^{-1}M)=a^{-1}a=1$. Fibers are disjoint and exhaust GL.
This establishes the count without assuming the first isomorphism
theorem. If that theorem is familiar, the quotient is on a backup slide.
}

\teachslide{In dimension two}{2}{
\[
\begin{aligned}
|\SL_2(\F_p)|
&=\frac{(p^2-1)(p^2-p)}{p-1}\\
&=\key{p(p^2-1)}.
\end{aligned}
\]
Lagrange's theorem: the order of a group element divides the group size.
\spacer
\key{Next time: apply this count to Fibonacci.}
}{
Factor $p^2-p=p(p-1)$ before cancelling. Briefly recall the element-order
definition: the least positive exponent giving the identity. Lagrange
applies to the cyclic subgroup generated by an element. This assumes
the course's earlier group material rather than reproving Lagrange.
}

\teachslide{Exit question}{2}{
For a field $K$ with $q$ elements, which group have we proved is cyclic?

What is its order?

Why does this not say that every nonzero element generates it?
}{
Expected answers: $K^\times$, with multiplication, of order $q-1$.
The proof finds an element of full order; many other elements can have
smaller orders, including $1$ when $q>2$. Ask students also to recall
the size $p(p^2-1)$ of SL(2,p). End at minute 50.
}

\backups

\backupslide{The determinant quotient}{
The determinant is surjective, with kernel $\SL_d(\F_p)$.
\[
\key{\GL_d(\F_p)/\SL_d(\F_p)\cong\F_p^\times.}
\]
A coset is determined by the common determinant of its matrices.

Multiplication of cosets corresponds to multiplication of determinants.
}{
This is the first isomorphism theorem. Kernel subgroups are normal,
so the quotient is defined. The explicit fiber bijection in the main
deck already proved the group-size formula without needing quotient
terminology. This optional slide answers the literal GL/SL question.
}
\backupslide{A different object with a cyclic unit group}{
For odd primes $p$ and $n\ge1$,
\[
\key{(\Z/p^n\Z)^\times\text{ is cyclic of order }p^{n-1}(p-1).}
\]
For $n>1$, the ring is not a field: $p\cdot p^{n-1}=0$.
\[
\begin{array}{c|c|c}
&\F_{p^n}&\Z/p^n\Z\\\hline
\text{Units}&p^n-1&p^{n-1}(p-1)\\
\text{Characteristic}&p&p^n
\end{array}
\]
}{
This is additional material, not an application of the finite-field
theorem to a residue ring. Units are exactly the classes not divisible
by $p$, giving $p^n-p^{n-1}$. Both factors in the displayed zero product
are nonzero when $n>1$. The next slides give a primitive-root lifting
proof. The field comparison assumes a field of order $p^n$; we have
not proved existence for every prime power in the main lectures.
}

\backupslide{Two arithmetics with nine elements}{
In $\Z/9\Z$, the powers of $2$ are
\[
1,2,4,8,7,5,\qquad2^6=1.
\]
These are all six units; the nonzero residues $3,6$ are not units.
\spacer
In $\F_9$, all eight nonzero elements are units,
and a primitive element has order $8$.
}{
Calculate successive powers modulo $9$. The nonzero residue $3$
satisfies $3^2=0$. A field of nine elements can instead be constructed
by adjoining a square root of $2$ to $\F_3$, using Lecture 2's method.
This contrasts a six-element unit group with an eight-element
multiplicative field group, despite equal sizes of the underlying sets.
}

\backupslide{First lift a primitive root to modulo $p^2$}{
Let $p$ be odd, and let $g$ generate $\F_p^\times$.

Choose $h=g$ or $h=g+p$ so that
\[
\key{h^{p-1}\not\equiv1\pmod{p^2}.}
\]
This is possible because
\[
(g+p)^{p-1}\equiv g^{p-1}+p(p-1)g^{p-2}\pmod{p^2},
\]
and the added term is nonzero modulo $p^2$.
}{
Choose any integer representative $g$ prime to $p$. The two choices
have the same primitive reduction modulo $p$. The linear binomial
term is divisible by $p$ but not $p^2$; all remaining terms are
divisible by $p^2$. Thus both powers cannot equal 1 modulo $p^2$.
Write $h^{p-1}=1+pc$ with $p\nmid c$. It would be false to say
every primitive representative automatically lifts without adjustment.
}

\backupslide{Each further lift contributes a factor of $p$}{
For odd $p$, if $u=1+p^rc$ with $r\ge1$ and $p\nmid c$, then
\[
u^p=1+p^{r+1}c',\qquad p\nmid c'.
\]
The first binomial term is $p^{r+1}c$; all later terms
are divisible by $p^{r+2}$.
\spacer
Starting with $h^{p-1}=1+pc$ gives
\[
\key{v_p\bigl(h^{(p-1)p^k}-1\bigr)=k+1.}
\]
}{
Define $v_p(m)$ as the exponent of $p$ dividing the nonzero integer
$m$, if needed. For $2\le j<p$, the binomial coefficient contributes
$p$, so the $j$th term has valuation at least $1+rj\ge r+2$.
The last term has valuation $rp\ge r+2$ since $p\ge3$.
After division by $p^{r+1}$, the result is $c$ modulo $p$, so
it remains nonzero. Induction gives the displayed formula.
This is a direct special-case proof, not an invocation of a general
lifting-the-exponent lemma.
}

\backupslide{The lifted element has full order}{
Modulo $p^n$, let $d=\ord(h)$.
\[
p-1\mid d,\qquad d\mid(p-1)p^{n-1}.
\]
Thus $d=(p-1)p^j$ for some $0\le j\le n-1$.

The lifting calculation shows the least such $j$ is $n-1$.
\[
\key{d=(p-1)p^{n-1}.}
\]
So $h$ generates every unit modulo $p^n$.
}{
Reduction modulo $p$ forces $p-1$ to divide the order; Lagrange
in the unit group gives the upper divisibility. The valuation formula
says $h^{(p-1)p^j}\equiv1\pmod{p^n}$ exactly when $j+1\ge n$.
This proves the exact order, with a single adjusted integer $h$
working for every $n$. The case $n=1$ is the finite-field theorem.
}

\backupslide{The exception starts at eight}{
\[
(\Z/2\Z)^\times=\{1\},\qquad(\Z/4\Z)^\times=\{1,3\}.
\]
Both are cyclic. Modulo $8$, every unit squares to $1$:
\[
1^2=3^2=5^2=7^2=1\pmod8.
\]
There is no element of order $4$.

For $n\ge3$,
\[
\key{(\Z/2^n\Z)^\times\cong C_2\times C_{2^{n-2}}.}
\]
}{
$C_m$ denotes a cyclic group of order $m$. The direct-product
generators are $-1$ and $5$. Repeated squaring gives
$v_2(5^{2^k}-1)=k+2$, starting with $5-1=4$:
if $u=1+2^rc$, $r\ge2$, $c$ odd, then
$u^2-1=2^{r+1}c(1+2^{r-1}c)$ and the last factor is odd.
So 5 has order $2^{n-2}$, and its powers are exactly the residues
equal to 1 modulo 4. Every odd residue is in this subgroup or its
negative, and $-1$ is outside it. This proves the direct product.
Its largest element order is $2^{n-2}$, smaller than its size
$2^{n-1}$, so it is not cyclic. The modulo-8 example alone suffices
for a short classroom comparison.
}

\backupslide{Why the field proof does not apply to residue rings}{
In $\Z/8\Z$, the degree-two polynomial
\[
X^2-1
\]
has four roots: $1,3,5,7$.
\spacer
\key{The polynomial root bound required a field.}

An abelian unit group need not be cyclic.
}{
The finite-field proof chose an element failing a low-degree polynomial
equation. With zero divisors, cancellation fails in the factor-theorem
argument. Thus the residue-ring theorem needs an additional lifting
argument, and the theorem fails for powers of 2 starting at 8.
This returns to Lecture 1's discussion of zero divisors.
}



\clearpage
\setcounter{slide}{0}\setcounter{elapsed}{0}
\section*{Lecture 4: Explaining Fibonacci periods}
\teachcover{Explaining\\Fibonacci periods}{1}{
The field containing the roots controls the powers
\[
\pi(11)=10,\qquad\pi(7)=16.
\]
}{
The first three lectures supplied fields, coordinates, cyclic
multiplicative groups, and counts of GL and SL. Today returns to
Fibonacci with a full 50 minutes for the order argument and Frobenius.
The prime binomial identity and Fermat's theorem are the remaining
prerequisites. Keep the calculation break before the modulo-11 example.
}

\teachslide{The structures we now understand}{2}{
For a finite field $K$ with $q$ elements, $K^\times$ is cyclic of order $q-1$.
\[
|\SL_2(\F_p)|=p(p^2-1).
\]
For Fibonacci, with $A=\mat1110$,
\[
\det A=-1,\qquad A^2=A+I.
\]
\key{For odd $p$, apply Lagrange to $A^2\in\SL_2(\F_p)$.}
}{
Recall that a primitive element has order $q-1$; not every element
is primitive. Matrix groups need not be cyclic. The determinant
of $A$ is -1, so its square rather than $A$ itself lies in SL.
Our field-group theorem is about multiplication of nonzero scalars.
}

\teachslide{The Fibonacci period is the matrix order}{4}{
For $t\ge1$,
\[
A^t=\begin{pmatrix}F_{t+1}&F_t\\F_t&F_{t-1}\end{pmatrix}.
\]
Since $F_{t-1}=F_{t+1}-F_t$,
\[
A^t=I\iff(F_t,F_{t+1})=(0,1)\pmod p.
\]
\[
\key{\pi(p)=\ord(A).}
\]
}{
The displayed power formula follows by induction from the recurrence,
and is the students' existing matrix approach. The equivalence needs
both state coordinates, not merely $F_t=0$. If a student asks why an
orbit period should equal the whole matrix order, this identity is
the specific Fibonacci justification; that equality need not hold for
an arbitrary matrix and arbitrary starting vector.
}

\teachslide{A first bound from $\SL_2$}{2}{
Let $p$ be odd. If $A^t=I$, determinants give $(-1)^t=1$.

Thus $\pi(p)$ is even, and
\[
\ord(A^2)=\frac{\pi(p)}2.
\]
Lagrange gives
\[
\key{\pi(p)\mid2p(p^2-1).}
\]
Can the roots of $X^2-X-1$ give a sharper answer?
}{
For an element of even order $t$, the least positive $r$ with $A^{2r}=I$
is $t/2$, because $t$ must divide $2r$. This proves the order identity.
The broad group-size divisibility uses no splitting information.
Its purpose is to motivate exploiting the special matrix, rather
than presenting the group size as an exact Fibonacci period.
}

\teachslide{Choose a field containing the roots}{3}{
Assume $p\ne2,5$.
\begin{itemize}
\item If $5$ is square in $\F_p$, use $K=\F_p$.
\item Otherwise, use $K=\{a+bs:s^2=5\}$ over $\F_p$.
\end{itemize}
\[
\alpha=\frac{1+s}{2},\qquad\beta=\frac{1-s}{2}.
\]
\[
\alpha+\beta=1,\quad\alpha\beta=-1,\quad\alpha\ne\beta.
\]
}{
In the split case, $s$ is an existing square root, not a new symbol.
In the nonsplit case, the proof from Lecture 2 generalizes directly
from $7$ to any odd $p$: $a^2-5b^2$ cannot vanish for a nonzero pair.
The prime $5$ is excluded because the roots coincide, and $2$ because
division by $2$ is unavailable. Both roots are nonzero since their
product is $-1$.
}

\teachslide{Diagonalize in the field where it makes sense}{4}{
For either root $r$, since $r^2=r+1$,
\[
A\begin{pmatrix}r\\1\end{pmatrix}
=r\begin{pmatrix}r\\1\end{pmatrix}.
\]
\[
P=\begin{pmatrix}\alpha&\beta\\1&1\end{pmatrix},
\qquad \det P=\alpha-\beta\ne0.
\]
\[
P^{-1}AP=\operatorname{diag}(\alpha,\beta),\qquad
\key{A^t=I\iff\alpha^t=\beta^t=1.}
\]
}{
Compute the product $(r+1,r)^{\mathsf T}$ and use the root equation.
This explicitly proves the eigenvectors and invertibility of the
eigenvector matrix, so no general diagonalization theorem is needed.
Changing the field lets us find $P$; it does not change whether a
power of the original matrix equals $I$. Its entries still lie in
the embedded prime field. This also gives the lcm of the two root orders.
}

\teachslide{A period is a simultaneous order}{2}{
We have shown
\[
A^t=I\iff\alpha^t=\beta^t=1.
\]
If $\ord(\alpha)=r$ and $\ord(\beta)=s$, the return times are
the common multiples of $r$ and $s$.
\[
\key{\pi(p)=\operatorname{lcm}(r,s).}
\]
}{
An element's return times are exactly the multiples of its order,
by division with remainder. We need both eigenvalues to return to 1.
This lcm statement is about simultaneous return, not the order of
the product of arbitrary elements. The coprime-product lemma from
Lecture 3 had an additional coprimality hypothesis.
}

\teachslide{Case one: the roots are in $\F_p$}{3}{
If $5$ is a square in $\F_p$, both roots lie in $\F_p^\times$.

Fermat's theorem gives
\[
\alpha^{p-1}=\beta^{p-1}=1.
\]
Therefore
\[
\key{A^{p-1}=I,\qquad\pi(p)\mid p-1.}
\]
}{
Remind students that the roots are nonzero, so Fermat applies.
The conclusion is divisibility, not equality. The new bound is
much sharper than the broad SL group-size bound because we have
used the eigenvalues of this particular matrix.
}

\teachslide{Your turn: roots modulo eleven}{3}{
In $\F_{11}$, use $s=4$ and $2^{-1}=6$.
\begin{enumerate}
\item Compute $\alpha=(1+s)/2$ and $\beta=(1-s)/2$.
\item Compute $\alpha^2$ and $\alpha^5$ by repeated squaring.
\item Which proper divisors of $10$ could still be the period?
\end{enumerate}
\spacer
\key{We already know that the period divides $10$.}
}{
Give students a minute to calculate before discussing. Expected
roots are 8 and 4; $8^2=9$, $8^4=4$, $8^5=-1$. These rule out
both 2 and 5 as return times, and therefore all proper divisors of
10. The next page reveals the answers and exact period.
}

\teachslide{Modulo eleven: an exact period}{2}{
Choose $s=4$ in $\F_{11}$. Then
\[
\alpha=8,\qquad\beta=4,\qquad\pi(11)\mid10.
\]
But
\[
8^5=-1,\qquad8^2=9\ne1\pmod{11}.
\]
Every proper divisor of $10$ divides $5$ or $2$.
\[
\key{\pi(11)=10.}
\]
}{
Compute $2^{-1}=6$ modulo $11$, giving roots $5\cdot6=8$ and
$(-3)\cdot6=4$. To check the powers, $8^2=9$, $8^4=4$, and
$8^5=32=-1$. If $A^5$ or $A^2$ were $I$, the corresponding power
of $8$ would be $1$. We have ruled out every proper divisor of the
known multiple $10$. The order of the other root need not be computed.
}

\teachslide{A Fibonacci root need not be primitive}{3}{
Modulo $11$, $\alpha=8$ has order $10$ and generates $\F_{11}^\times$.

Modulo $7$, the roots live in $\F_{49}$, where
\[
|\F_{49}^\times|=48.
\]
But $\alpha=4+4s$ satisfies $\alpha^8=-1$ and $\alpha^{16}=1$.
\[
\key{\ord(\alpha)=16,\quad\text{so it is not a multiplicative generator}.}
\]
}{
This comparison connects cyclicity directly to the order problem.
In the field from Lecture 2, $\alpha^2=5+4s$, $\alpha^4=5s$,
and $\alpha^8=(5s)^2=125=-1$. Thus its exact order is 16.
Nevertheless $s=2\alpha-1$, so $\alpha$ does generate this field
over $\F_7$. Generating the extension field and generating its
multiplicative group are different requirements. Frobenius will
explain the order bound without needing these special calculations.
}

\teachslide{A new symmetry: Frobenius}{4}{
In characteristic $p$,
\[
\key{(u+v)^p=u^p+v^p.}
\]
The intermediate binomial coefficients are divisible by $p$.

Also $(uv)^p=u^pv^p$ and $1^p=1$.

On a finite field, $u\mapsto u^p$ is injective, hence bijective:
\[
u^p=v^p\quad\Longrightarrow\quad(u-v)^p=0\quad\Longrightarrow\quad u=v.
\]
}{
Recall why $p$ divides $\binom pj$ for $0<j<p$: the numerator
contains a factor $p$, while $j!$ is prime to $p$. The power map
preserves the field operations and is a bijection, so it is a field
automorphism, called Frobenius. A nonzero field element cannot have
a zero power. Finiteness is needed for the injective-to-surjective
step; no infinite-field claim is required.
}

\teachslide{Which elements does Frobenius fix?}{3}{
Every $a\in\F_p$ satisfies $a^p=a$.

The fixed elements are roots of
\[
X^p-X.
\]
A nonzero polynomial of degree $p$ has at most $p$ roots.

The prime field already supplies $p$ roots.
\[
\key{u^p=u\iff u\in\F_p.}
\]
}{
For nonzero prime-field elements use Fermat; zero is immediate.
Recall the polynomial root bound proved in Lecture 3. The polynomial is nonzero
since its leading coefficient is $1$, even in characteristic $p$.
This argument identifies the fixed elements in our extension without
assuming the desired conjugation rule.
}

\teachslide{Case two: Frobenius exchanges the roots}{3}{
Suppose $5$ is nonsquare in $\F_p$, with $s^2=5$ in $K$.
\[
(s^p)^2=5^p=5.
\]
Thus $s^p$ is $s$ or $-s$.

It cannot be $s$, because $s\notin\F_p$.
\[
\key{s^p=-s,\qquad\alpha^p=\beta,\qquad\beta^p=\alpha.}
\]
}{
The factorization $X^2-5=(X-s)(X+s)$ in $K$ and absence of zero
divisors give exactly those two square roots. They are distinct since
$p$ is odd and $s\ne0$. If $s^p=s$, the preceding fixed-element result
would put $s$ in the prime field, contradicting nonsquareness.
Frobenius fixes $1$ and $2^{-1}$, so applying it to the formulas for
$\alpha,\beta$ interchanges them.
}

\teachslide{The minus sign tells us the return time}{3}{
\[
\alpha^{p+1}=\alpha\beta=-1,
\qquad\beta^{p+1}=-1.
\]
Diagonalization gives
\[
\key{A^{p+1}=-I,\qquad A^{2(p+1)}=I.}
\]
Therefore
\[
\pi(p)\mid2(p+1),\qquad\pi(p)\nmid p+1.
\]
At time $p+1$, $(F_{p+1},F_{p+2})=(0,-1)$.
}{
Conjugating the scalar matrix $-I$ leaves it unchanged. Since the
characteristic is odd, $-I\ne I$. Applying the matrix to $(1,0)$
gives $(F_{p+2},F_{p+1})=(-1,0)$, hence the displayed state in the
usual order. This connects the minus sign to the earlier warning
that a Fibonacci zero alone is not a restart.
}

\teachslide{Modulo seven: the other exact period}{3}{
Since $5$ is nonsquare modulo $7$,
\[
A^8=-I,\qquad A^{16}=I.
\]
Thus $\pi(7)\mid16$ but $\pi(7)\nmid8$.

Every proper divisor of $16$ divides $8$.
\[
\key{\pi(7)=16.}
\]
}{
This exact-order argument avoids a long power calculation in the
extension. Ask students why $A^8\ne I$ rules out every proper divisor
of $16$. For a direct Frobenius check in the field constructed last
time, $s^7=s(s^2)^3=125s=-s$. The general proof explains why this
calculation is not a coincidence.
}

\teachslide{What the algebra explained}{2}{
For primes $p\ne2,5$,
\[
\begin{array}{c|c}
5\text{ is square in }\F_p&\pi(p)\mid p-1\\\hline
5\text{ is nonsquare in }\F_p&\pi(p)\mid2(p+1)
\end{array}
\]
\begin{itemize}
\item Group size supplies a broad order bound.
\item The field containing the roots supplies a sharper bound.
\item Further power checks identify the least period.
\end{itemize}
}{
Recap the distinction between existence of a period, a known multiple,
and the least period. The matrix construction proves existence for
all moduli; group structure and the root field give increasingly
specific information for prime moduli. Prime $29$ is an optional
example where the bound is not attained.
}

\teachslide{Exit question}{3}{
Suppose $p$ is odd and $5$ is nonsquare modulo $p$.

A student argues: ``$F_{p+1}=0$, so the period is $p+1$.''

What is missing, and what does the matrix actually tell us?
\spacer
\key{Use both entries of the Fibonacci state.}
}{
Expected answer: a restart requires $(F_{p+1},F_{p+2})=(0,1)$,
but the state is $(0,-1)$ because $A^{p+1}=-I$. The period divides
$2(p+1)$ and does not divide $p+1$; it need not equal $2(p+1)$.
End the main lecture here at minute 50. Use optional slides only
if time remains or as material for a subsequent meeting.
}

\backups
\backupslide{Why the prime five is different}{
Modulo $5$, $X^2-X-1=(X-3)^2$.
\[
A=3I+N,\qquad N=\mat3112,\qquad N^2=0.
\]
For $t\ge1$,
\[
A^t=3^tI+t3^{t-1}N.
\]
$A^t=I$ requires $5\mid t$ and $4\mid t$.
\[
\key{\pi(5)=20.}
\]
}{
Compute $A-3I$ modulo $5$ to get the displayed nonzero nilpotent matrix.
The off-diagonal entry in $A^t$ forces $t=0$ in the field, hence
$5\mid t$, because $3$ is nonzero. Then $A^t=3^tI$ and $3$ has
order $4$ modulo $5$: $3^2=-1$ and $3^4=1$. These two divisibilities
are also sufficient, giving least time $20$. The repeated root
produces a nilpotent term instead of two independent eigenvectors.
}

\backupslide{The prime two}{
Modulo $2$, the successive Fibonacci states are
\[
(0,1)\ \longmapsto\ (1,1)\ \longmapsto\ (1,0)
\ \longmapsto\ (0,1).
\]
\[
\key{\pi(2)=3.}
\]
Here $-1=1$, and division by $2$ is unavailable.
}{
The short state cycle gives the exact answer immediately. This also
explains why the odd-characteristic parity argument and the root
formula with denominator $2$ were explicitly restricted.
}

\backupslide{A bound need not be the least period}{
Modulo $29$, $11^2=5$, so $\pi(29)\mid28$.

The roots are $\alpha=6$ and $\beta=24$.
\[
6^7=-1,\quad6^2=7\ne1;
\qquad24^7=1,\quad24\ne1.
\]
Their orders are $14$ and $7$.
\[
\key{\pi(29)=\operatorname{lcm}(14,7)=14.}
\]
}{
Check $6^7=-1$ by repeated squaring: $6^2=7$, $6^4=20$,
$6^7=20\cdot7\cdot6=28$. Thus $6$ has order $14$, since every
proper divisor of $14$ divides $7$ or $2$. The other root is
$-6^{-1}$, so $24^7=(-1)^7(6^7)^{-1}=1$; since $7$ is prime
and $24\ne1$, its order is $7$. This is a concrete proper-divisor
example, not just a caution about the word divides.
}

\backupslide{An exact-period procedure}{
Start with a proven multiple $B$ of $\pi(p)$.
\begin{enumerate}
\item Set $t=B$ and factor $t$.
\item For each prime divisor $q$ of $t$, test $A^{t/q}=I$.
\item If the test succeeds, replace $t$ by $t/q$ and repeat.
\end{enumerate}
When no prime factor can be removed, $t=\pi(p)$.
}{
Matrix repeated squaring performs the tests efficiently. At every
stage the true order divides $t$. If $t$ were still larger than
the order, some prime factor of $t/\pi(p)$ could be removed while
leaving a multiple of the order. Thus the stopping rule is exact.
For $29$, start at $28$ and remove one factor $2$ to reach $14$.
}

\clearpage
\section*{Course connections and scope}
The primary materials are the local number-theory course book:
\texttt{chapters/modular\_fibonacci.tex} for reversible states and period
definitions, \texttt{chapters/finite\_fields.tex} for quadratic extension
arithmetic, Frobenius, and Fibonacci root orders, and
\texttt{chapters/matrix\_algebra.tex} for the matrix toolkit.
The group and binomial chapters supply the prerequisites for Lectures 3 and 4.
The ring and vector-space development and the ordered-basis count expand
that route for these lectures.

The lectures prove that every finite field has prime-power order. They
construct the quadratic fields required when $5$ is nonsquare over an
odd prime field. They prove cyclicity of finite-field multiplicative groups, using
Lagrange, the polynomial root bound, and a coprime-order product lemma.
They do not prove existence and uniqueness of finite fields of every
prime-power size or quadratic reciprocity. Those classification results
are not needed for the displayed Fibonacci divisibilities. The exact-period examples use
only powers and known multiples of the order.
\end{document}
